\expandafter\def\csname HistoricalNoteText26\endcsname{Proprie quidem ex analysi praecedenti tantummodo sequitur, $-G'$ et $-G''$ satisfacere debere aequationi 13, unde dubium esse videtur, annon liceat, utramque $-G'$ et $-G''$ \textit{eidem} radici
negativae aequalem ponere, prorsus neglecta radice tertia. Sed facile perspicietur, siquidem aequationis radix secunda et tertia sint inaequales, ex $-G'=\allowbreak -G''$ sequi $\gamma'=\gamma''$, $\alpha'=\alpha''$, $\beta'=\beta''$, et proin
$-\alpha'\alpha''-\beta'\beta''+\allowbreak \gamma'\gamma''=\allowbreak -\alpha'\alpha'-\beta'\beta'+\allowbreak \gamma'\gamma'=1$, quod aequationi quartae in [1] est contrarium. Conf. quae
infra de casu duarum radicum aequalium aequationis 13 dicentur.}
